% Generated draft: requires mathematical and source review.
\documentclass[11pt]{article}
\usepackage[margin=25mm]{geometry}
\usepackage{fontspec}
\setmainfont{Latin Modern Roman}
\usepackage{amsmath,amssymb,mathtools,unicode-math}
\setmathfont{Latin Modern Math}
\usepackage{enumitem,needspace,xurl,hyperref}
\setlength{\parindent}{0pt}
\setlength{\parskip}{5pt}
\setlength{\emergencystretch}{3em}
\newcommand{\N}{\mathbb{N}}
\newcommand{\Z}{\mathbb{Z}}
\newcommand{\Q}{\mathbb{Q}}
\newcommand{\R}{\mathbb{R}}
\newcommand{\C}{\mathbb{C}}
\newcommand{\F}{\mathbb{F}}
\newcommand{\A}{\mathbb{A}}
\newcommand{\PP}{\mathbb{P}}
\newcommand{\E}{\mathbb{E}}
\newcommand{\eps}{\varepsilon}
\newcommand{\dd}{\,\mathrm{d}}
\begin{document}
\section*{UnsolvedMath ID: 30004322}
\textbf{Source problem code:} catalogue entry
\subsection{Definitions and mathematical statement}
The supplied question concerns the singular points of a line configuration in $\PP^2$. The base field and the meaning of ``line configuration'' are not specified. The following is a conditional formulation using the usual interpretation as a finite reduced union of distinct projective lines over an algebraically closed field $K$; neither this field convention nor any restriction on its characteristic is confirmed by the record.

Let $X=\PP^2_K$, let $\mathcal A=\bigcup_{j=1}^{d}L_j$ be such a union, and let $S=\{x_1,\ldots,x_r\}$ be its set of distinct singular points, namely the points incident with at least two component lines. Assume $r\geq 1$, as needed for the expression below. Write $H=\mathcal O_X(1)$. Define the multipoint Seshadri constant by
\[
\varepsilon(H;x_1,\ldots,x_r)
=
\inf_{\substack{C\subset X\ \text{integral curve}\\ \sum_{i=1}^{r}\operatorname{mult}_{x_i}(C)>0}}
\frac{\deg C}{\sum_{i=1}^{r}\operatorname{mult}_{x_i}(C)},
\]
where $\operatorname{mult}_{x_i}(C)=0$ when $x_i\notin C$.

Reading ``on one line'' as allowing every projective line in $X$, define
\[
k=\max_{\ell\subset X\ \text{projective line}}\#(S\cap\ell).
\]
Under these conventions, the supplied question asks whether, for every such configuration,
\[
\varepsilon(\mathcal O_{\PP^2_K}(1);x_1,\ldots,x_r)=\frac{1}{k}.
\]
The record does not clarify whether the maximum defining $k$ is instead restricted to the component lines $L_1,\ldots,L_d$. That alternative gives a different formulation and remains unresolved here. Configurations with no singular points are also not addressed by the source wording, since then $k=0$ and $1/k$ is undefined.

The supplied background and literature assessment discuss global-generation and very-ampleness thresholds on a surface, rather than the displayed line-configuration question. Their intended connection to this question is not supplied, so they cannot resolve its missing hypotheses or establish its current mathematical status.
\subsection{Short English statement}
For the singular points of a finite arrangement of projective lines, the question asks whether the multipoint Seshadri constant of $\mathcal O_{\PP^2}(1)$ equals the reciprocal of the largest number of these points on a line. The base field, admissible arrangements, and whether that line must belong to the arrangement require clarification. The supplied literature assessment concerns a different formulation.
\subsection{Sources}
\begin{itemize}
\item Local UnsolvedMath registry, \texttt{data/unsolvedmath/problems.json}, supplied record with permanent ID \texttt{30004322}, title ``Seshadri Constants of Line Arrangement Singularities,'' exact source code \texttt{catalogue entry}. No individual proposer is supplied. The record gives proposed year 2019 and labels the problem \texttt{open} and the statement \texttt{corrected\_verified}; these are supplied registry labels, not independent verification. Its verification note says that only formula delimiters, notation, grammar, and typography were repaired, with mathematical content and scope preserved.
\item Supplied report attribution: \emph{Mini-Workshop: Seshadri Constants (2020)}, Oberwolfach Reports, DOI 10.4171/owr/2019/53, pp.~3271--3313. \url{https://doi.org/10.4171/owr/2019/53}. The background describes Oberwolfach Reports as documenting mathematical workshops and their open-problem discussions. The supplied citation uses 2020, while the DOI, proposed year, and year tag refer to 2019; their relationship has not been verified.
\item Supplied literature source: Ein and Lazarsfeld, \emph{Seshadri constants on smooth surfaces}. \url{https://www.math.stonybrook.edu/~roblaz/Reprints/Ein.Laz.Seshadri.Consts.Smooth.Sfs.pdf}.
\item Supplied literature source label: ``Oberwolfach Reports, global generation problem.'' \url{https://doi.org/10.4171/owr/2019/53}.
\item Supplied literature assessment checked at \texttt{2026-08-22}: the record labels the status open and states that Kollar-type constructions give unbounded very-ampleness thresholds on one surface, but do not settle the extracted globally-generated threshold question, nor did subsequent work located in that assessment. This is supplied dated curation derived from the OWR material, not independently verified. The supplied background expressly requests independent verification of time-sensitive status and cited claims. This assessment concerns a threshold question whose connection to the present line-configuration statement is unspecified.
\end{itemize}
\end{document}
